\section{Introduction}
Robots operating over extended periods encounter new manipulation tasks
while continuing to rely on previously acquired skills. Continual imitation
learning updates a shared policy from sequentially available demonstrations.
New-task updates also change representations and decisions used by
earlier tasks, making retention central to long-term robot learning.

Experience replay preserves historical demonstrations for subsequent
training. Multimodal latent replay reduces the storage burden through
compact representations, while incremental feature adjustment regularizes
task representations \citep{yu2026mlr}. M2Distill further preserves
multimodal features and action distributions across consecutive stages
\citep{roy2025m2distill}. These methods establish replay and behavioral
regularization as complementary tools for continual learning.

Two connected questions motivate this work. First, which demonstration
experience should remain available under a limited replay quota?
Uniform random selection does not explicitly optimize representativeness,
while purely geometric selection does not impose retention of a random
reference subset. We combine these selection principles through protected
random anchors and representative coverage within the same quota.

Second, which histories should supply behavior-retention targets?
Demonstrations record expert-controlled interaction. Successful autonomous
executions record histories generated by the learned controller and
environmental feedback. Constraints evaluated only on demonstration
histories do not directly constrain outputs on additional autonomous
histories. Retaining these histories connects task completion experience
to the training of later tasks.

We introduce a framework combining RA-Coreset and SHBR. RA-Coreset repairs
a representative subset to include protected random anchors. SHBR maintains
a separate success-history memory with behavior references fixed at task
acquisition. It adjusts retention using trajectory support and selects
candidate updates using current-task and retention probes. Demonstrations
continue to drive learning and correction, while successful histories supply
additional behavioral references for subsequent learning.

Our contributions are:
\begin{enumerate}
\item RA-Coreset, a fixed-quota replay selector combining protected random
anchors with representative demonstration coverage.
\item SHBR, a success-history retention mechanism combining fixed execution
references, trajectory-support weighting, and controlled update selection.
\end{enumerate}
\draftnote{Complete citation-supported motivation for replay selection.
Do not claim that prior methods lack policy distillation.}
